\acknowledgement

时光荏苒，转眼就到本科毕业的时候了。这篇论文的完成，离不开许多人的帮助，在这里一一感谢。

首先要感谢我的指导老师。从选题到方案设计，从实验调试到论文撰写，老师在每个阶段都给了我很多具体的指导。老师的严谨和对问题的判断力，我一直很佩服，也跟着学了不少。遇到瓶颈的时候，老师总能几句话帮我理清思路，绕开死胡同。

感谢实验室的同学们。项目开发期间，大家经常凑在一起讨论方案，互相看代码、提意见，这种氛围让很多问题提前暴露出来。特别感谢几位同学在系统架构和代码实现上给我的建议，你们的想法让我看到了一些自己完全没想到的地方。

也感谢班级的同学们。大家研究方向各不相同，平时聊起来反而更容易碰撞出意想不到的想法。谢谢你们在生活上的关照，一起熬夜、一起吃饭的日子，我会记得。

感谢我的家人。父母这么多年一直在背后支持我，从来不给压力，只是在我需要的时候说几句实在话。累的时候、怀疑自己的时候，想到他们就又能撑下去了。

回头看整个毕设，从对知识图谱和自然语言处理只有模糊概念，到能独立搭出一套完整的方案并跑通，这个过程中我慢慢学会了两件事：遇到不懂的就去学，一个人搞不定就找人帮忙。这两点听上去简单，做起来并不容易。

最后，衷心感谢在百忙之中审阅本论文的各位专家和老师，谢谢你们提出的宝贵意见。

谨以此文，献给所有在我大学四年学习和生活中给予帮助的人们。

\vskip 2ex
\rightline{华中科技大学}
\rightline{计算机科学与技术学院}
\rightline{2026年6月}
